\section*{The Council Debate}\stepcounter{section}\phantomsection\addcontentsline{toc}{section}{The Council Debate}
{\entryfont The gathering quickly reveals that there is no common vision for what should follow the closing of the portals. Although opinions throughout Ty Coeden differ considerably, three broad positions dominate the discussion. These are not formal factions, and many Caelthir sympathize with more than one of them. The following viewpoints provide the DM with the principal arguments that can be raised throughout the debate.}
\subsection*{We Must Remain Hidden}
{\entryfont For some, the disappearance of the portals changes little about Ty Coeden's relationship with the outside world. Isolation has protected the Caelthir for generations, allowing them to preserve their customs and safeguard the strange magic of the Woods of Lomond without interference from rulers, merchants, or opportunists beyond the forest.

Those supporting this view fear that greater contact may eventually bring unwanted attention to Ty Coeden. Fey magic, unusual creatures, rare plants, and knowledge accumulated over generations could all attract those who would seek to exploit them. They advocate maintaining the village's secrecy, continuing regular patrols, and limiting contact with outsiders to exceptional circumstances.}
\begin{DndReadAloud}
	The portals are gone, but the world beyond these woods has not changed with them. Kings still wage wars, merchants still covet what is rare, and strangers still seek what they do not understand. Ty Coeden endured because it remained hidden. We should not surrender that safety merely because one danger has ended.
\end{DndReadAloud}
\eject
\subsection*{Our Purpose Is Over}
{\entryfont Others see the closing of the portals primarily as the end of a burden carried for generations. Scouts watched the forest for new rifts, healers prepared for frightened arrivals, and countless Caelthir shaped their lives around protecting the woods from whatever might emerge next. That responsibility is now gone.

Those holding this view do not necessarily oppose contact with the outside world, but neither do they believe the Caelthir must immediately replace one obligation with another. For the first time in generations, individuals may be free to decide what they want their lives to become. They advocate reducing the constant vigilance surrounding the former portals and allowing Ty Coeden time to discover its future gradually.}
\begin{DndReadAloud}
	For generations, we lived waiting for the next portal. Scouts watched the woods, healers prepared for frightened arrivals, and warriors trained for whatever might emerge. Yesterday, that burden ended. Perhaps we do not need another duty to replace it. For once, let us ask not what is required of us, but what we want.
\end{DndReadAloud}
\subsection*{We Belong to Two Worlds}
{\entryfont The most outward-looking voices argue that the Caelthir can no longer treat their life in Fife as temporary. Many among them were born beneath the trees of the Woods of Lomond and have never seen the Feywild at all. Their families, memories, and dead are rooted in the Material Plane, regardless of where their ancestors once came from.

Those supporting this view believe that Ty Coeden should cautiously begin establishing relationships with the people beyond the forest. Trade, shared knowledge, and peaceful contact need not mean abandoning Fey heritage. Instead, they see an opportunity for the Caelthir to become part of Fife while remaining distinctly themselves.}
\begin{DndReadAloud}
	We speak of the Feywild as home, yet many among us have never seen it. We were born beneath these trees, learned these paths, and buried our dead in this soil. That does not make us less Caelthir. We can preserve what we were while accepting what we have become. Perhaps Fife is our home as well.
\end{DndReadAloud}
\subsection*{Running the Debate}
{\entryfont The three positions should serve as points of tension rather than a sequence of prepared speeches. Allow speakers to interrupt, challenge, question, and occasionally agree with one another. A cautious elder may accept that Fife has become home while still opposing open contact, just as an advocate for integration may insist that Ty Coeden must retain its independence.

No position is intended to be objectively correct, and the debate should remain unresolved. Its purpose is to expose the uncertainty within Ty Coeden and create natural opportunities for the player characters to enter the discussion before the council eventually settles upon a cautious compromise.}
\eject